% SOURCE_AUTHORITY: sga5_fr_workpass.tex, lines 10410--10463 (LF-normalized unit).
% SOURCE_SHA256: 791F4EFFC5E02832D5D77ED03518C8156D6F07E4C8238B03545DB93D883FBB28
\begin{statement}{Corolario 7.2}
Sean $X$ un $k$-esquema proyectivo liso, conexo, de dimensión $n$, $Y$ un subesquema de $X$ que es intersección de $d$ secciones hiperplanas de $X$, y $L$ un $\mathbb Z_\ell$-haz constante torcido constructible $(\ell\ne\operatorname{car}(k))$. Entonces:
\begin{enumerate}[label=(\roman*)]
\item el homomorfismo de imagen inversa
\[
\alpha_i:H^i(X,L)\to H^i(Y,L|Y)
\]
es un isomorfismo para $i\le n-d-1$, y un monomorfismo para $i=n-d$. Si, además, $L$ es localmente libre constructible, el conúcleo de $\alpha_{n-d}$ es libre de torsión.
\item supongamos que $Y$ es liso sobre $k$, de dimensión $n-d$. Entonces el homomorfismo de Gysin
\[
\beta_i:H^i(Y,(L|Y)(-d))\to H^{i+2d}(X,L)
\]
es un isomorfismo para $i\ge n-d+1$ y un epimorfismo para $i=n-d$. Si, además, $L$ es localmente libre constructible, el epimorfismo $\beta_{n-d}$ es escindido y su núcleo es libre de torsión.
\end{enumerate}
\end{statement}

\emph{Demostración.} Demostremos (i); la afirmación (ii) se deduce por dualidad de Poincaré. En el caso general, basta un simple paso al límite a partir de 7.1 (i). Cuando $L$ es un $\mathbb Z_\ell$-haz localmente libre, se aplica el lema siguiente al morfismo de imagen inversa
\[
R\Gamma(X,L)\to R\Gamma(Y,L|Y)
\]
de $D(\mathbb Z_\ell)$; la hipótesis (iii) se verifica gracias a 7.1 (i).

\begin{statement}{Lema 7.2.1}
Sean $V$ un anillo de valoración discreta, $k_V$ su cuerpo residual y $u:K\to L$ un morfismo de $D(V)$. Se supone que $K$ y $L$ tienen cohomología de tipo finito, y se denota por $M$ el cilindro de la aplicación $u$. Entonces, para todo $m\in\mathbb Z$, las condiciones siguientes son equivalentes:
\begin{enumerate}[label=(\roman*)]
\item $H^p(M)=0$ para $p\le m-1$, y $H^m(M)$ es libre de torsión.
\item para todo $V$-módulo $N$, el morfismo
\[
H^p(u\otimes^L\operatorname{id}_N):H^p(K\otimes^L N)\to H^p(L\otimes^L N)
\]
es un isomorfismo para $p\le m-1$ y un monomorfismo para $p=m$;
\item[(ii bis)] para todo $V$-módulo $N$, se tiene
\[
H^p(M\otimes^L N)=0\qquad(p\le m-1);
\]
\item el morfismo
\[
H^p(u\otimes^L\operatorname{id}):H^p(K\otimes^L k_V)\to H^p(L\otimes^L k_V)
\]
es un isomorfismo para $p\le m-1$ y un monomorfismo para $p=m$;
\item[(iii bis)] se tiene
\[
H^p(M\otimes^L k_V)=0\qquad(p\le m-1).
\]
\end{enumerate}
Además, estas condiciones equivalentes implican:
\begin{enumerate}[label=(\roman* bis)]
\item el morfismo
\[
H^p(u):H^p(K)\to H^p(L)
\]
es un isomorfismo para $p\le m-1$, y un monomorfismo de conúcleo libre para $p=m$.
\end{enumerate}
\end{statement}
